\documentclass[conference]{IEEEtran}
\usepackage{graphicx}
\usepackage{subcaption}
\usepackage{float}
\usepackage{amsmath}
\usepackage{amssymb}
\usepackage{algorithm}
\usepackage{algpseudocode}
\usepackage{booktabs}
\usepackage{multirow}

\usepackage[colorlinks=true, urlcolor=cyan, linkcolor=blue, citecolor=blue]{hyperref}

\title{RBE549: Project 3 --- Einstein Vision\\
\large Tesla-Inspired 3D Scene Understanding and Visualization}

\author{
\IEEEauthorblockN{Nishanth Adithya Chandramouli}
\IEEEauthorblockA{Robotics Engineering\\
Worcester Polytechnic Institute\\
nchandramouli@wpi.edu}
\and
\IEEEauthorblockN{Chaitanya Sriram Gaddipati}
\IEEEauthorblockA{Robotics Engineering\\
Worcester Polytechnic Institute\\
csgaddipati@wpi.edu}
}

\begin{document}
\maketitle

\begin{abstract}
We present \emph{Einstein Vision}, a Tesla FSD-inspired 3D visualization
pipeline for autonomous driving scenes.  Given monocular video from a front-
facing camera on a 2023 Tesla Model S, the system detects and tracks
vehicles, pedestrians, lane markings, traffic lights, road signs, and
miscellaneous scene objects, then renders a photorealistic Blender scene
from a bird's-eye perspective.  The pipeline spans three phases of
increasing sophistication: Phase~1 covers detection, depth estimation, and
3D localisation; Phase~2 adds multi-object tracking, pedestrian pose, and
fine-grained classification; Phase~3 introduces brake/turn-signal
detection, ego-motion-compensated motion classification, and
collision-risk prediction.  Our key architectural contributions are (1)
fusing YOLO11~\cite{ultralytics2024yolo11} and YOLO-World~\cite{cheng2024yoloworld}
to cover both standard COCO classes and open-vocabulary infrastructure
objects in a single NMS-unified pass; (2) using CLIP~\cite{radford2021clip}
for vehicle sub-type classification and speed-limit OCR fallback; and
(3) applying a RAFT~\cite{teed2020raft} optical flow field together with
fundamental-matrix RANSAC and Sampson-distance epipolar tests to separate
independently-moving vehicles from camera ego-motion.
\end{abstract}

% ============================================================
\section{Introduction}
% ============================================================

Autonomous vehicles must build a rich, real-time model of their surroundings
from sensors that are inherently noisy and partial.  Modern driver-assistance
HMIs---most visibly Tesla's Full Self-Driving visualisation dashboard---
distil this model into a synthetic third-person view that conveys depth,
motion, and semantic meaning at a glance.  Replicating this pipeline from
scratch is a demanding integration task: it requires accurate single-frame
detection, temporally consistent tracking, metric depth estimation,
kinematic state estimation, and high-fidelity 3D rendering, all running
in a coherent data flow.

Our system, \emph{Einstein Vision}, addresses this challenge.  The input is
raw MP4 video from the front camera of a 2023 Tesla Model S driving through
13 diverse urban scenes.  The output is a Blender-rendered video that mirrors
the Tesla FSD aesthetic: road surfaces, lane markings as B\'ezier curves,
vehicle and pedestrian meshes placed at metric 3D positions with correct yaw,
traffic signals with correct colour and arrow states, and Phase-3 overlays for
brake lights, turn signals, motion arrows, and collision-risk highlights.

The pipeline is implemented as a two-stage process.  The \emph{perception
stage} (Python, PyTorch) processes each sampled frame to produce a structured
\texttt{FrameData} record (detections, lanes, depth, flow) serialised to
MessagePack and NumPy archives.  The \emph{rendering stage} (Blender/Python,
\texttt{bpy}) reads those archives headlessly and renders each frame to a PNG.
Both stages are driven by a single YAML configuration file so any model can be
swapped without changing pipeline code.

% ============================================================
\section{Phase 1: Core Perception}
% ============================================================

\subsection{Object Detection}

We fuse two YOLO-family models to cover the full set of required
driving-scene classes.

\textbf{YOLO11~\cite{ultralytics2024yolo11}} (weights: \texttt{yolo11x.pt})
handles the 80 standard COCO classes.  We discard non-driving classes
(furniture, food, sports equipment) and keep only:
\emph{person, bicycle, car, motorcycle, bus, truck, traffic light, stop sign,
fire hydrant}.

\textbf{YOLO-World~\cite{cheng2024yoloworld}} (weights:
\texttt{yolov8x-worldv2.pt}) handles open-vocabulary classes that COCO does
not cover: \emph{traffic cone, traffic cylinder, traffic pole, dustbin,
barrel, speed limit sign, road arrow marking, speed bump}.  We set a
custom class list at runtime via \texttt{model.set\_classes()}, which
inserts a text-embedding branch without retraining.

The two models run independently on each frame.  Their detection lists are
concatenated and then de-duplicated with a single batched NMS pass
(IoU~$=~0.45$) so genuinely overlapping boxes from the same class compete
against each other while boxes for different classes do not suppress one
another.  Both models share the same confidence threshold of $0.30$.

\subsection{Vehicle Sub-classification and Speed-Limit OCR}

\textbf{Vehicle sub-types.}
YOLO11 reports all light passenger vehicles as \texttt{car}.  To distinguish
sedan / SUV / hatchback / pickup we crop each \texttt{car} bounding box and
run a CLIP~\cite{radford2021clip} ViT-B/32 similarity against four text
prompts, one per sub-type.  The text features are pre-encoded at
initialisation so the per-frame cost is a single image-encoder forward pass
plus a dot product.

\textbf{Speed limit signs.}
Detected \texttt{speed limit sign} crops are first tested for aspect
ratio~$\le~2.0$ to filter diamond-shaped warning signs.  Valid crops are
upscaled to a minimum 160~px short-side, converted to grayscale, Otsu-
thresholded on the upper two-thirds (where the number lives), and passed
to EasyOCR with a digit allowlist.  If EasyOCR returns a digit string that
matches a valid US speed limit ($\{15, 20, \ldots, 75\}$ in steps of 5),
that value is recorded.  When OCR fails (blurry crop, severe occlusion), the
same crop is classified with CLIP against fifteen speed-limit text prompts as
a fallback.

\subsection{Lane Detection}

We use a fine-tuned Mask~R-CNN~\cite{he2017maskrcnn} (ResNet-50-FPN backbone)
with 7 output classes: background plus six lane types (solid-white, dashed-
white, solid-yellow, double, and two additional generic classes mapped to the
nearest canonical type).  The model outputs per-instance soft masks at full
input resolution.

Each binary mask is thresholded at $0.50$ and requires at least 150 active
pixels to be accepted.  From the surviving pixels we fit a second-degree
polynomial $x = f(y)$ to obtain a lane curve aligned with the driving
direction, then uniformly sample 10 B\'ezier control points along the curve.
These 2D pixel-space points are later lifted to 3D ego-frame coordinates
using the ground-plane Inverse Perspective Mapping (IPM) approach described
in Section~\ref{sec:depth}.

A lightweight IoU-based lane tracker assigns consistent lane IDs across
frames by matching each new detection to the nearest existing track by
centroid proximity (threshold 80~px).  Unmatched tracks are aged out after
10 frames.

\subsection{Traffic Light Classification}

\textbf{Colour.}
For each \texttt{traffic light} bounding box we divide the crop into three
equal horizontal strips (top~=~red, middle~=~yellow, bottom~=~green) and
compute the mean brightness of the HSV $V$-channel in each strip.  The strip
with the highest mean brightness above a minimum threshold ($V > 60$) is
taken as the active colour.  This 3-strip heuristic is fast, requires no
colour calibration, and handles different traffic-light housing styles that
HSV hue-range matching misclassifies under non-neutral lighting.

\textbf{Arrow direction.}
After the colour strip is identified we attempt arrow detection on that strip.
The colour-specific HSV mask is morphologically cleaned (open then close with
a $3\times3$ elliptical kernel) and the largest contour is extracted.  A
circularity test
\begin{equation}
  C = \frac{4\pi \cdot \text{area}}{\text{perimeter}^2}
\end{equation}
rejects round bulbs ($C > 0.55$).  For non-circular contours the image
moments give a central axis orientation:
\begin{equation}
  \theta = \tfrac{1}{2}\arctan\!\left(\frac{2\mu_{11}}{\mu_{20} - \mu_{02}}\right)
\end{equation}
Near-vertical axes ($|\theta| > 55^\circ$) are labelled
\emph{straight}; near-horizontal axes are resolved to \emph{left} or
\emph{right} by comparing the centroid $x$-position to the medial axis
of the contour points.

\subsection{Depth Estimation and 3D Localisation}
\label{sec:depth}

\textbf{Monocular depth.}
We use Depth Anything V2 Metric Outdoor (Large)~\cite{li2024depthanything},
accessed via the HuggingFace Transformers interface.  The model produces
float32 metric depth maps (metres) at native frame resolution via bicubic
upsampling of the internal feature-resolution prediction.  Batches of frames
are processed together on the GPU to amortise model startup overhead.

\textbf{Lifting detections to 3D.}
For each detection the depth is sampled as the median of a $9\times9$ pixel
patch centred on the bounding-box centroid.  Patch-median is preferred over
a single-pixel lookup because depth models are often inaccurate at object
boundaries.  If the sampled depth is below $1.0$~m (indicating an unreliable
prediction) we fall back to a perspective-height estimate:
\begin{equation}
  d = \frac{f_y \cdot H_{\text{real}}}{h_{\text{px}}}
\end{equation}
where $H_{\text{real}}$ is a class-specific real-world height (e.g.\
$1.5$~m for a car) and $h_{\text{px}}$ is the bounding-box height in pixels.
The 3D camera-frame point is then:
\begin{equation}
  \mathbf{x}_{\text{cam}} =
  \begin{bmatrix}
    \dfrac{(u - c_x)\,d}{f_x} \\[6pt]
    \dfrac{(v - c_y)\,d}{f_y} \\[6pt]
    d
  \end{bmatrix}
\end{equation}
and is transformed into the ego (vehicle body) frame via the camera extrinsic
matrix.

\textbf{Lane lifting.}
Lane B\'ezier control points are lifted to 3D via IPM: given a flat-ground
assumption with camera height $h_c$ above ground and pitch $\phi$, each
pixel $(u, v)$ maps to a ground-plane point.  When a dense depth map is
available the IPM ground-plane depth is cross-validated against the depth
map; the depth-map value takes precedence if the two agree within a factor
of 2.

\textbf{Yaw estimation.}
For each vehicle we estimate heading angle geometrically.  A vehicle of
length $L$ and width $W$ viewed at yaw $\theta$ subtends a projected ground
width proportional to $W\cos\theta + L\sin\theta$.  We invert this relationship
using the measured projected width (bounding-box width $\times$ depth / $f_x$)
to recover $\theta \in [0, \pi/2]$, then resolve the front/rear ambiguity
from the sign of the horizontal position relative to the ego vehicle.

% ============================================================
\section{Phase 2: Advanced Features}
% ============================================================

\subsection{Multi-Object Tracking with DeepSORT}

We use DeepSORT~\cite{wojke2017deepsort} (via \texttt{deep-sort-realtime})
with a MobileNet appearance embedder.  Detections from the fused YOLO
pipeline are forwarded to the tracker each frame; the tracker returns
confirmed track IDs after $n_{\text{init}} = 3$ consecutive detections.
Lost tracks are retained for up to 30 frames before deletion, which
handles short occlusions without creating ghost detections.

Consistent track IDs are critical for Phase-3 features: brake-light
detection uses a per-track temporal buffer to distinguish single-frame
reflections from sustained illumination, and the collision predictor
maintains a position history per track to estimate velocity.

\subsection{Pedestrian Pose Estimation}

Pedestrian pose is estimated with YOLO11x-pose, which predicts the 17
standard COCO body keypoints (nose, eyes, ears, shoulders, elbows, wrists,
hips, knees, ankles) together with per-keypoint confidence scores from the
bounding-box crop.  For schema compatibility with the 133-keypoint RTMW
format used elsewhere, the 17-keypoint output is zero-padded: the first 17
entries hold the COCO keypoints and the remaining 116 (face, hand, foot
keypoints) are set to zero.  The Blender pose renderer reads the
non-zero entries and skips the zero-padded skeleton segments, so the
rendering degrades gracefully to a partial skeleton rather than
drawing artefact limbs.

\subsection{Traffic Sign Detection and OCR}

Stop signs and speed limit signs are detected by YOLO-World and YOLO11
respectively (see Section II-A).  After detection, EasyOCR extracts the
numeric value from speed limit sign crops as described in Section II-B.
Confirmed detections are serialised with their 3D ego-frame position so
the Blender renderer can spawn the correct sign asset at the correct location.

\subsection{Road Arrow Detection}

Ground arrow markings (turn arrows, keep-right, etc.) are detected by
YOLO-World with the label \texttt{road arrow marking}.  The bounding-box mask
is extracted and the largest contour found.  Principal Component Analysis on
the contour points yields two eigenvectors; the primary eigenvector's angle
relative to the image vertical is thresholded to classify the arrow as
\emph{straight}, \emph{left}, or \emph{right}.

\subsection{Miscellaneous Object Detection}

Traffic cones, traffic cylinders, traffic poles, dustbins, and barrels are
detected by YOLO-World using the open-vocabulary class list in
\texttt{configs/models.yaml}.  Each object class has a corresponding 3D
asset in the Blender asset library.  Because these objects have small or
irregular shapes, the depth-map fallback (perspective height) is used for
objects shorter than 1~m.

% ============================================================
\section{Phase 3: Bells and Whistles}
% ============================================================

\subsection{Brake Light and Turn Signal Detection}

We analyse three regions of interest (ROIs) extracted from each tracked
vehicle's bounding box.

\textbf{Brake lights.}
The bottom $40\,\%$ of the bounding box (the tail-light band) is converted
to HSV and the fraction of pixels satisfying a bright-red HSV range
($H \in [0^\circ,10^\circ]\cup[160^\circ,179^\circ]$, $S > 80$,
$V > 120$) is computed.  A brake-light-on event is triggered when this
red-pixel density exceeds $5\,\%$.

\textbf{Turn signals.}
The left and right $30\,\%$ vertical strips of the lower $60\,\%$ of the
bounding box are analysed for amber ($H \in [8^\circ,28^\circ]$, $S > 110$,
$V > 120$) pixel density.  Turn signals blink, so a temporal history of 8
frames is maintained per track.  A turn signal is classified as active when
(a) the mean amber density over the history exceeds $0.5\times$ the threshold
($2\,\%$) \emph{and} (b) at least one frame in the recent history exceeds the
full threshold ($4\,\%$).  This two-condition gate rejects brief reflections
(which fail condition~a) and persistent but faint ambient colour (which fail
condition~b).

\textbf{Blender rendering.}
Active brake lights drive the emissive intensity of the tail-light material
slots on the vehicle mesh.  Left or right turn-signal activation drives
the corresponding indicator emissive slot.  The light animator
(\texttt{blender/light\_animator.py}) interpolates emissive values across
frames so blinking is visible in the rendered video.

\subsection{Moving vs.\ Parked Vehicle Classification}

Classifying independently-moving vehicles in the presence of strong
ego-motion is a fundamentally ill-posed problem for pure optical-flow
magnitude thresholding.  We address this with a two-stage approach.

\textbf{Stage 1 — background flow estimation.}
We compute the median optical flow vector $(b_x, b_y)$ over the entire frame.
Because moving objects typically occupy less than $50\,\%$ of the image, the
frame-wide median is a robust estimate of the camera-induced background flow.
The per-vehicle residual flow is then
\begin{equation}
  \Delta\mathbf{u}_i = \mathbf{u}_i - (b_x, b_y)
\end{equation}
where $\mathbf{u}_i$ is the median flow inside the vehicle's bounding box.

\textbf{Stage 2 — Sampson distance test.}
We additionally estimate the fundamental matrix $F$ of the current frame
transition from 800 randomly sampled pixel correspondences using
RANSAC~\cite{hartley2004multiple} (\texttt{cv2.FM\_RANSAC}).  $F$ encodes
the camera's rotational and translational ego-motion.  For a stationary
scene point the Sampson distance
\begin{equation}
  d_S(\mathbf{x}, \mathbf{x}') =
    \frac{(\mathbf{x}'^T F \mathbf{x})^2}
         {(F\mathbf{x})_0^2 + (F\mathbf{x})_1^2 + (F^T\mathbf{x}')_0^2 + (F^T\mathbf{x}')_1^2}
\end{equation}
should be near zero; for an independently-moving point it will be large.
A vehicle is classified as \emph{moving} when \emph{either}:
\begin{itemize}
  \item its residual flow magnitude $\|\Delta\mathbf{u}_i\| > 2.0$~px/frame, \emph{or}
  \item its Sampson distance $d_S > 4.0$.
\end{itemize}
The dual-criterion design handles the failure modes noted in previous work:
the flow threshold catches vehicles moving faster than the ego vehicle
(visible residual) while the Sampson test catches vehicles moving at similar
speed to ego (small residual but significant epipolar deviation).

Optical flow itself is computed by RAFT (Large, \texttt{raft-things}
weights)~\cite{teed2020raft} with 20 recurrent update iterations, operating
on pairs of consecutive sampled frames after divisibility-by-8 padding.

\subsection{Collision Risk Prediction (Extra Credit)}

The collision predictor maintains a rolling 10-frame position history per
confirmed track.  Each call to \texttt{update()} appends the current
ego-frame 3D position $(x, y, z)$ to the per-track deque and fits a linear
velocity:
\begin{equation}
  \mathbf{v}_i = \frac{1}{N-1}\sum_{k=1}^{N-1}
    \bigl(\mathbf{p}_{i,k} - \mathbf{p}_{i,k-1}\bigr)
\end{equation}
Tracks with fewer than 3 entries are skipped to avoid false positives
from newly-appeared objects.

Predicted positions 15 frames ahead ($\approx 2$~s at the default
7.2~fps sampling rate) are computed as:
\begin{equation}
  \hat{\mathbf{p}}_i = \mathbf{p}_i + \mathbf{v}_i \cdot T_{\text{pred}}
\end{equation}
All pedestrian--vehicle track pairs are checked: if the predicted
3D Euclidean distance between a pedestrian and a vehicle falls below
$5.0$~m, both tracks are flagged with \texttt{collision\_risk = True}.
The Blender renderer responds by applying a red emissive overlay to the
flagged meshes, giving the operator an immediate visual warning.

% ============================================================
\section{Pipeline Architecture}
% ============================================================

\subsection{Data Flow}

Figure~\ref{fig:pipeline} illustrates the end-to-end data flow.  The
perception stage (\texttt{scripts/run\_perception.py}) reads frames at a
configurable stride, runs all models in sequence, and writes one
MessagePack file (\texttt{FrameData}) and one or two NumPy archives (depth,
flow) per frame.  The rendering stage (\texttt{blender/scene\_builder.py})
reads those archives headlessly under Blender's embedded Python interpreter
and renders each frame to a PNG that is later assembled into a video.

\begin{figure}[H]
  \centering
  \includegraphics[width=\linewidth]{figs/pipeline_diagram.png}
  \caption{End-to-end pipeline architecture.  Perception (left) produces
  serialised per-frame records; Blender rendering (right) consumes them
  headlessly.}
  \label{fig:pipeline}
\end{figure}

\subsection{Per-Frame Processing Order}

Within the perception stage, models run in the following order to respect
data dependencies:
\begin{enumerate}
  \item \textbf{Depth} — Depth Anything V2 produces the metric depth map.
  \item \textbf{Detection} — YOLO11 + YOLO-World fused with cross-model NMS.
  \item \textbf{Sub-classification} — CLIP vehicle sub-type and OCR speed limit.
  \item \textbf{Tracking} — DeepSORT assigns/updates track IDs.
  \item \textbf{3D Lifting} — \texttt{lift\_detections()} fills \texttt{pos\_3d} and \texttt{depth}.
  \item \textbf{Yaw} — Geometric yaw estimation per vehicle.
  \item \textbf{Traffic lights} — HSV 3-strip colour + arrow detection.
  \item \textbf{Lanes} — Mask R-CNN segmentation + IoU lane tracker.
  \item \textbf{Pose} — YOLO11x-pose for person detections.
  \item \textbf{Phase 3} — Brake/turn (BrakeTurnDetector), flow (RAFT),
    motion (MotionClassifier), collision (CollisionPredictor).
\end{enumerate}

\subsection{Configuration and Model Swapping}

All model choices, weight paths, confidence thresholds, and hyperparameters
live in \texttt{configs/models.yaml}.  The factory module
(\texttt{src/perception/factory.py}) reads the \texttt{model} key under each
section and instantiates the corresponding concrete subclass of the relevant
abstract base class.  Swapping, say, the depth estimator from
\texttt{depth\_anything\_v2} to \texttt{depth\_pro} requires a single YAML
edit and no code changes.

Table~\ref{tab:config} summarises the configuration used for all results
presented in this report.

\begin{table}[H]
\centering
\caption{Perception pipeline configuration.}
\label{tab:config}
\begin{tabular}{lp{4.2cm}}
\toprule
Component & Configuration \\
\midrule
Depth model        & Depth Anything V2 Metric Outdoor-L \\
Detection          & YOLO11x + YOLO-World v8x (conf $0.30$, IoU $0.45$) \\
Sub-classifier     & CLIP ViT-B/32 + EasyOCR \\
Lane detector      & Mask R-CNN ResNet-50-FPN (conf $0.50$) \\
Tracker            & DeepSORT MobileNet ($\max\_\text{age}=30$, $n_\text{init}=3$) \\
Pose               & YOLO11x-pose (17 COCO keypoints) \\
Optical flow       & RAFT-Large, \texttt{raft-things}, 20 iters \\
Brake threshold    & Red density $> 5\%$ \\
Turn threshold     & Amber density $> 4\%$, history 8 frames \\
Motion: flow       & Residual $> 2.0$~px/frame \\
Motion: Sampson    & $d_S > 4.0$ \\
Collision horizon  & 15 frames, warn $< 5.0$~m \\
\bottomrule
\end{tabular}
\end{table}

% ============================================================
\section{Blender Rendering}
% ============================================================

The Blender rendering stage reads the serialised \texttt{FrameData} records
and rebuilds the 3D scene for each frame from scratch.

\textbf{Asset placement.}  The asset manager (\texttt{blender/asset\_manager.py})
maintains a library of \texttt{.blend} mesh files, one per object class
(sedan, SUV, pickup, truck, bus, motorcycle, bicycle, pedestrian, traffic
light pole-and-head assembly, stop sign, speed limit sign, traffic cone,
traffic cylinder, dustbin).  For each frame, vehicle and pedestrian meshes
are instantiated at their 3D ego-frame positions with yaw applied as a
Z-axis rotation.  Per-track colour assignment ensures that the same vehicle
carries the same mesh colour throughout the video, aiding visual continuity.

\textbf{Lane curves.}  Lanes are rendered as B\'ezier curves via Geometry
Nodes (\texttt{blender/lane\_renderer.py}).  The 3D B\'ezier control points
from the perception stage are fed directly to a curve object, and a Geometry
Nodes modifier extrudes the curve profile into a flat lane-marking strip.
Solid and dashed lane types use different material slots with animated
dash patterns.

\textbf{Lighting.}  A dynamic HDRI sky plus a sun lamp tracks a simulated
solar arc through the sequence.  The light animator
(\texttt{blender/light\_animator.py}) adjusts vehicle tail-light and
indicator emissive intensities per frame based on the brake and turn-signal
flags from Phase~3.

\textbf{Camera.}  The Blender camera is configured from the calibration YAML
($f_x, f_y, c_x, c_y$, height above ground, pitch) so that the virtual
camera matches the real sensor geometry.

% ============================================================
\section{Qualitative Results}
% ============================================================

Figures~\ref{fig:phase1_out}, \ref{fig:phase2_out}, and \ref{fig:phase3_out}
show representative rendered frames from Phase~1, Phase~2, and Phase~3
respectively.

\begin{figure}[H]
  \centering
  \includegraphics[width=\linewidth]{figs/phase1_render.png}
  \caption{Phase~1 output: vehicles and pedestrians placed at metric 3D
  positions, lane B\'ezier curves rendered as white/yellow markings, stop
  sign and traffic light assets placed at detected locations.}
  \label{fig:phase1_out}
\end{figure}

\begin{figure}[H]
  \centering
  \includegraphics[width=\linewidth]{figs/phase2_render.png}
  \caption{Phase~2 output: pedestrian skeletons (17 COCO keypoints), vehicle
  sub-type labels (sedan / SUV / pickup), speed limit sign with OCR-extracted
  value, and road arrow markings.}
  \label{fig:phase2_out}
\end{figure}

\begin{figure}[H]
  \centering
  \includegraphics[width=\linewidth]{figs/phase3_render.png}
  \caption{Phase~3 output: brake light emissive glow (red) on a braking
  sedan, turn-signal amber glow on a turning vehicle, motion arrows on
  independently-moving vehicles, and red collision-risk overlays on a
  pedestrian and nearby vehicle.}
  \label{fig:phase3_out}
\end{figure}

Figure~\ref{fig:debug_overlay} shows the 2D debug overlay (generated by
\texttt{scripts/make\_debug\_video.py}) that is produced alongside the
Blender render.  This overlay draws bounding boxes, track IDs, depth
estimates, lane masks, and optical flow magnitude on the original camera
image, providing a diagnostic view of the perception stage outputs.

\begin{figure}[H]
  \centering
  \includegraphics[width=\linewidth]{figs/debug_overlay.png}
  \caption{2D debug overlay showing detection bounding boxes with class
  labels and track IDs, lane mask, and depth values at each bbox centre.
  Green~=~moving, grey~=~parked.}
  \label{fig:debug_overlay}
\end{figure}

% ============================================================
\section{Challenges and Failure Modes}
% ============================================================

\textbf{Yaw estimation.}
The geometric yaw estimator assumes each vehicle subtends a clean
rectangular silhouette.  Partially occluded vehicles, tight front-quarter
views, and vehicles at the frame boundary all violate this assumption,
leading to yaw errors of $20$--$40^\circ$.  A learning-based 3D bounding
box estimator (e.g.\ DD3D or BEVFusion) would reduce this error but
requires calibrated stereo or LiDAR ground truth for training, which is
not available for this dataset.

\textbf{Depth scale drift.}
While Depth Anything V2 is metric, the absolute scale drifts slightly
across scenes.  Objects at ranges beyond approximately 40~m show $10$--$20\,\%$
depth errors, causing distant vehicles to appear too large or too close in
the Blender render.  A two-point scale calibration using known lane-width
measurements would improve consistency.

\textbf{Traffic light colour under night conditions.}
The 3-strip brightness method assumes the active bulb is the brightest strip.
In night scenes with strong ambient light pollution (sodium vapour street
lamps produce a yellow-orange glow), the brightness baseline is elevated,
occasionally causing misclassification of red lights as yellow or green.
A learned detector or per-scene white-balance calibration would address this.

\textbf{Ego-motion compensation.}
The RANSAC fundamental-matrix approach requires sufficient parallax across
the frame.  On near-straight highway sections where the forward translation
dominates and there is little rotation, the estimated $F$ is degenerate and
the Sampson test reverts to the residual-flow threshold alone.  A homography-
based approach that models pure translation may be more robust in these cases.

\textbf{Lane detection on faded markings.}
The Mask R-CNN lane detector was fine-tuned on a dataset of clearly marked
lanes.  Several scenes contain heavily worn or snow-covered lane markings that
the model misses entirely.  Pseudo-label training on the target domain or a
complementary drivable-area segmentation model could bridge this gap.

% ============================================================
\section{Author Contributions}
% ============================================================

\textbf{Nishanth Adithya Chandramouli:}
Phase-3 motion classification (RAFT integration, Sampson distance, ego-motion
compensation), collision predictor, brake/turn detector, Blender rendering
pipeline (scene\_builder, light\_animator, lane\_renderer), debug visualisation.

\textbf{Chaitanya Sriram Gaddipati:}
Object detection pipeline (YOLO11 + YOLO-World fusion, NMS), vehicle
sub-classifier (CLIP + EasyOCR), traffic light HSV classification and arrow
detection, depth integration (Depth Anything V2), 3D localisation and yaw
estimation, DeepSORT tracking integration, report writing.

% ============================================================
\begin{thebibliography}{10}

\bibitem{ultralytics2024yolo11}
Ultralytics,
``YOLOv11: A state-of-the-art real-time object detection model,''
2024.
\url{https://docs.ultralytics.com/models/yolo11/}

\bibitem{cheng2024yoloworld}
T.~Cheng, L.~Song, Y.~Ge, W.~Liu, X.~Wang, and Y.~Shan,
``YOLO-World: Real-Time Open-Vocabulary Object Detection,''
\emph{arXiv:2401.17270}, 2024.

\bibitem{radford2021clip}
A.~Radford \emph{et al.},
``Learning Transferable Visual Models From Natural Language Supervision,''
in \emph{Proc.\ ICML}, 2021.

\bibitem{he2017maskrcnn}
K.~He, G.~Gkioxari, P.~Doll{\'a}r, and R.~Girshick,
``Mask R-CNN,''
in \emph{Proc.\ ICCV}, 2017.

\bibitem{li2024depthanything}
Y.~Li \emph{et al.},
``Depth Anything V2,''
\emph{arXiv:2406.09414}, 2024.

\bibitem{teed2020raft}
Z.~Teed and J.~Deng,
``RAFT: Recurrent All-Pairs Field Transforms for Optical Flow,''
in \emph{Proc.\ ECCV}, 2020.

\bibitem{wojke2017deepsort}
N.~Wojke, A.~Bewley, and D.~Paulus,
``Simple Online and Realtime Tracking with a Deep Association Metric,''
in \emph{Proc.\ ICIP}, 2017.

\bibitem{hartley2004multiple}
R.~Hartley and A.~Zisserman,
\emph{Multiple View Geometry in Computer Vision}, 2nd~ed.
Cambridge University Press, 2004.

\bibitem{mousavian20163dbb}
A.~Mousavian, D.~Anguelov, J.~Flynn, and J.~Ko{\v{s}}eck{\'a},
``3D Bounding Box Estimation Using Deep Learning and Geometry,''
in \emph{Proc.\ CVPR}, 2017.

\bibitem{zhou2022detic}
X.~Zhou, V.~Koltun, and P.~Kr{\"a}henb{\"u}hl,
``Detecting Twenty-Thousand Classes using Image-level Supervision,''
in \emph{Proc.\ ECCV}, 2022.

\end{thebibliography}

\end{document}
